\section{Einleitung}

Der Einfluss der Menschheit auf die Natur und diesen Planeten steigt immer weiter an.
(Quelle?)
Im gleichen Zuge schreitet die Automatisierung und Digitalisierung voran.
(Quelle?)
Viele Gemeinden, Kommunen und Einzelpersonen suchen daher nach der Verknüpfung von Technik und Umwelt, um mit den Daten aus der Umwelt die Natur zu erhalten, und ein zusammenleben mit ihr zu ermöglichen. (Green Planet Quelle?)

Durch viele Umweltkatastrophen wird dieser Standpunkt noch verstärkt.

\subsection{Motivation}

Die Häufigkeit von Waldbränden nimmt stetig zu, der Faktor Mensch ist seiner Umgebung in vielerlei Hinsicht nicht bewusst.
In den meisten Flüssen, welche Großstädte durchfließen herrscht Badeverbot und keine Trinkwasserqualität.

Die eigene Umgebung wieder zu renaturalisieren und Flächen zu schaffen, welche mit Artenvielfalt und gesunder Natur aufwarten können ist eine große motivation.

Dafür müssen Menschen sich ihrer Umgebung bewusst werden, und sich mit ihr beschäftigen.

Ein weg dahin sind sogenannte "Green Technologies", Begrünung in der Stadt, und Integration von Technik mit Natur.
Gewächshäuser die Automatisch die Temperatur regeln, Alarmsysteme im Wald die bei Dürre ausschlagen um Waldbrände zu verhindern, oder Bojen die konstant die Wasserqualität überwachen, um den Maritimen Ökosystemen eine Stimme in unserer Gesellschaft zu geben.





\subsection{Stand der Technik}

Im Bereich der Maker-Community und IoT gibt es viele Projekte welche sich mit Agrarsensoren und Automatisierung beschäftigen.

Verteilte Systeme wie LoRaWan und auch Anwendungen über Mobilfunk sind zu finden.

Ein Vorbildprojekt, welches zu großen Teilen in dieser BA weitergeführt wird, ist die Vision der Spreeboje von Jakob Kukula. Dabei geht es um eine schwimmende Plattform zur Wasserqualitätsuntersuchung mit dem Ziel, Flüsse und Seen auf der Welt zu verknüpfen und den Menschen den Zustand ihrer Aquatischen Umgebung näherzubringen.

Ein Proof of Concept von dem Projekt der Waldboje von Sara Reichert und Jakob Kukula lag dieser BA als Ansatz zugrunde.

\subsection{Ziel der Arbeit}

Das Resultat dieser Arbeit soll eine Plattform für die eigenständige Datenmessung und Weiterleitung in das Internet sein.
Es soll autonom, über einen langen Zeitraum, in Urbanen sowie Städtischen Gebieten arbeiten.
Es soll mit verschiedenen Sensoren ausrüstbar sein, und eine Basis für weitere "Green Tech" Arbeiten darstellen.
Eine eigene Stromversorgung, sowie Energiemanagement sind vorgesehen.
Im Bezug auf die Spreeboje sollen auch mögliche Aktuatoren wie Licht-Feedback evaluiert werden.
Die Baupläne so wie Anleitungen zum Nachbau sollen verfügbar gemacht, und DIY-Kommunen auf der Welt angeregt werden, eigene Module zu bauen.
Die Vision ist eine Globale Verknüpfung der Module um Aufmerksamkeit auf die Umgebung und die Natur zu lenken.


% System mit dem Umweltdaten gesammelt und zu einem geschickt werden

% Konkrete Vision, eine Plattform aus günstigen und verfügbaren Teilen.
% Für Maker-Communitys auf der Welt.
% Mit Schnittstellen \ac{GPIO} und \ac{ADC} für genügen Sensoren und Aktuatoren
% Zum einfachen Nachbauen

%  >> Reproduzierbares System, Tutorial, für Maker, um Umweltdaten zu sammeln, bereitzustellen und sich auf der Welt zu verknüpfen
%     >> Vereinfachung / Baukit für leute die nicht so viel Hardware Ahnung haben


% Ziel: das Befinden \& den Zustand der Natur greifbarer für Menschen zu machen.

\subsection{Abgrenzung}

Das entwickelte System soll nicht für die Komerzielle nutzung und nicht als Konkurrenz hochentwickelter schon vorhandener technischer Lösungen zur Erkennung von Waldbränden, Frühwarnsystemen, oder Wasserüberwachungsanlagen dienen.
Es soll als ein Beispielprojekt die Verbindung von Mensch und Natur darstellen.

Der Fokus der Arbeit soll auch nicht auf den Sensoren zur Datenaufnahme liegen, sondern auf der Plattform für deren Verarbeitung und Weiterleitung, sowie der Energieversorgung.

% TODO: Mehr sagen dass die Datenaufnahme AKA Sensoren nicht mein Tea ist.



% Andere BA:
% -Einleitung (Was ist in den Letzten Jahren Passiert? Womit beschäftigt sich diese Arbeit)
%     >> Viele Waldbrände, Generelle Umweltverschmutzung, Automatisierung der Umwelt, und des Lebens IoT.
%     -Motivation (So nen bisschen Grund, Was soll das ganze bringen, wo sind vllt Defizite in unserer Gesellschaft die damit gefüllt werden)
%     >> Das System kann Maker-Communitys ansprechen und Menschen mehr mit ihrer Umwelt verbinden
%     -Stand der Technik (Was gibt es schon für Arbeiten in die Richtung)
%     >> Arduino, ESP32, Adafruit viele DIY-Projekte für Pflanzenüberwachung, Gewächshausautomatisierung, 
%     -Ziel der Arbeit (Entwicklung Platine die Position bestimmen kann im Orbit)
%     >> Reproduzierbares System, Tutorial, für Maker, um Umweltdaten zu sammeln, bereitzustellen und sich auf der Welt zu verknüpfen
%     >> Vereinfachung / Baukit für leute die nicht so viel Hardware Ahnung haben
%     -Abgrenzung (Kein endgültiges perfektes Produkt)
%     >> Kein Komerzielles Produkt, eher ein Proof of concept, und ein Startpunkt um für Projekte zu dienen
%     -Struktur und Methodik der Arbeit (Vorstellung wo in der BA was erläutert wird)
%     >> Design Process suchen? (Systems engineering Struktur)
%     >> Vorstellung der Kapitel (in Kapitel 1 wird das und das vorgestellt, Rahmenbedingungen werden in kap. 2 Vorgestellt etc.
    

% \subsection{Struktur und Methodik}




% \section{Grundlagen?}


% \subsection{Mikroprozessor}

% \subsection{Solarzellen}

% \subsection{}

% \subsection{Datenverlauf}
